% Paper F: Differentiable Constraint Solving via Semiring Relaxation
% "The Learning Paper" - miniKanren 1-Bit Paper Suite
%
% For God so loved the world that he gave his only begotten Son,
% that whoever believes in him should not perish but have eternal life.
% John 3:16

\documentclass[11pt,a4paper]{article}

% Packages
\usepackage[utf8]{inputenc}
\usepackage{amsmath,amssymb,amsthm}
\usepackage{algorithm}
\usepackage{algpseudocode}
\usepackage{booktabs}
\usepackage{graphicx}
\usepackage{hyperref}
\usepackage{listings}
\usepackage{xcolor}
\usepackage{xspace}
\usepackage{tikz}
\usepackage[normalem]{ulem}  % For \sout (strikethrough)
\usetikzlibrary{shapes,arrows,positioning,fit,calc}

% Theorem environments
\newtheorem{theorem}{Theorem}
\newtheorem{lemma}[theorem]{Lemma}
\newtheorem{definition}[theorem]{Definition}
\newtheorem{proposition}[theorem]{Proposition}
\newtheorem{corollary}[theorem]{Corollary}

% Code listings
\lstset{
  basicstyle=\ttfamily\small,
  keywordstyle=\color{blue},
  commentstyle=\color{gray},
  stringstyle=\color{red},
  showstringspaces=false,
  breaklines=true
}

% Include shared definitions
% Shared Definitions for miniKanren 1-Bit Paper Suite
% Include with: % Shared Definitions for miniKanren 1-Bit Paper Suite
% Include with: % Shared Definitions for miniKanren 1-Bit Paper Suite
% Include with: \input{../shared_chirho/shared_defs_chirho}
%
% For God so loved the world that he gave his only begotten Son,
% that whoever believes in him should not perish but have eternal life.
% John 3:16

% ============================================================================
% Core Definitions (shared across all papers)
% ============================================================================

% The Finite-Domain miniKanren Kernel
\newcommand{\fdmk}{FD-miniKanren\xspace}  % Finite-Domain miniKanren

% ============================================================================
% Scope Clarification (include in each paper's introduction)
% ============================================================================

\newcommand{\scopeclarification}{%
\textbf{Scope Clarification.}
This paper addresses the \emph{Finite-Domain miniKanren Kernel} (\fdmk),
where variable domains are finite bitmasks.
Standard miniKanren unification over unbounded term structures uses a separate
reference implementation; the bit-parallel approach accelerates the constraint
propagation core, not full structural unification.
}

% ============================================================================
% Common Notation
% ============================================================================

% Domains
\newcommand{\dom}[1]{D_{#1}}          % Domain of variable x
\newcommand{\domand}{\land}            % Domain intersection (AND)
\newcommand{\domor}{\lor}              % Domain union (OR)

% Search
\newcommand{\state}{(\vec{D}, \sigma)} % Search state
\newcommand{\findChirho}{\text{find}}      % Union-find find
\newcommand{\unionChirho}{\text{union}}    % Union-find union

% Tensors
\newcommand{\tensor}[1]{T_{#1}}        % Tensor for relation R
\newcommand{\contract}{\circledast}    % Tensor contraction

% Semirings
\newcommand{\softand}{\odot}           % Soft AND (multiply)
\newcommand{\softor}{\oplus}           % Soft OR (prob sum)

% ============================================================================
% Hierarchical Domain Scaling
% ============================================================================

% Domain type naming (semantic macros for consistency across papers)
% Using semantic names to avoid confusion with Rust type identifiers
\newcommand{\domainSmallChirho}{\texttt{BitVec64Chirho}}
\newcommand{\domainMediumChirho}{\texttt{Hierarchical4kChirho}}
\newcommand{\domainLargeChirho}{\texttt{Hierarchical256kChirho}}
\newcommand{\domainDiffChirho}{\texttt{DiffHierarchical4kChirho}}

% Measured performance values (from Rust benchmarks, Jan 2026)
\newcommand{\domainSmallTimeChirho}{420 ps}
\newcommand{\domainMediumTimeChirho}{39 ns}
\newcommand{\domainLargeTimeChirho}{2.5 $\mu$s}
\newcommand{\domainLargeOpsChirho}{400K ops/sec}

% ============================================================================
% Paper Suite Reference
% ============================================================================

% Use these to cite sibling papers
\newcommand{\paperA}{Paper A: Bit-Parallel Finite-Domain Relational Search}
\newcommand{\paperB}{Paper B: Relations as Sparse Boolean Tensors}
\newcommand{\paperC}{Paper C: Bridging Infinite Domains with Hash-Consed Term IDs}
\newcommand{\paperD}{Paper D: Tabling and Mode-Driven Goal Scheduling}
\newcommand{\paperE}{Paper E: Fully Fused Logic Accelerator}
\newcommand{\paperF}{Paper F: Differentiable Constraint Solving via Semiring Relaxation}

% ============================================================================
% Acknowledgments Template
% ============================================================================

\newcommand{\standardacknowledgments}{%
Above all, I thank my Lord and Savior Jesus Christ for saving a wretch like me.

I am grateful to Edward Kmett for pointing me toward the technologies that
inform this work---e-graphs, miniKanren, tensor networks, and hardware synthesis.

I also thank my father, Roger Cowan, for his patient support.

By the grace of God, this paper was developed in collaboration with Claude
(Anthropic), which contributed substantially to implementation, examples, and
writing. Google Gemini and OpenAI's GPT-5.2 Codex provided valuable feedback
on drafts.
}

% ============================================================================
% Example environment (for papers that use it)
% ============================================================================
\newenvironment{exampleChirho}{\par\noindent\textbf{Example.}\itshape}{\par}

% Soli Deo Gloria

%
% For God so loved the world that he gave his only begotten Son,
% that whoever believes in him should not perish but have eternal life.
% John 3:16

% ============================================================================
% Core Definitions (shared across all papers)
% ============================================================================

% The Finite-Domain miniKanren Kernel
\newcommand{\fdmk}{FD-miniKanren\xspace}  % Finite-Domain miniKanren

% ============================================================================
% Scope Clarification (include in each paper's introduction)
% ============================================================================

\newcommand{\scopeclarification}{%
\textbf{Scope Clarification.}
This paper addresses the \emph{Finite-Domain miniKanren Kernel} (\fdmk),
where variable domains are finite bitmasks.
Standard miniKanren unification over unbounded term structures uses a separate
reference implementation; the bit-parallel approach accelerates the constraint
propagation core, not full structural unification.
}

% ============================================================================
% Common Notation
% ============================================================================

% Domains
\newcommand{\dom}[1]{D_{#1}}          % Domain of variable x
\newcommand{\domand}{\land}            % Domain intersection (AND)
\newcommand{\domor}{\lor}              % Domain union (OR)

% Search
\newcommand{\state}{(\vec{D}, \sigma)} % Search state
\newcommand{\findChirho}{\text{find}}      % Union-find find
\newcommand{\unionChirho}{\text{union}}    % Union-find union

% Tensors
\newcommand{\tensor}[1]{T_{#1}}        % Tensor for relation R
\newcommand{\contract}{\circledast}    % Tensor contraction

% Semirings
\newcommand{\softand}{\odot}           % Soft AND (multiply)
\newcommand{\softor}{\oplus}           % Soft OR (prob sum)

% ============================================================================
% Hierarchical Domain Scaling
% ============================================================================

% Domain type naming (semantic macros for consistency across papers)
% Using semantic names to avoid confusion with Rust type identifiers
\newcommand{\domainSmallChirho}{\texttt{BitVec64Chirho}}
\newcommand{\domainMediumChirho}{\texttt{Hierarchical4kChirho}}
\newcommand{\domainLargeChirho}{\texttt{Hierarchical256kChirho}}
\newcommand{\domainDiffChirho}{\texttt{DiffHierarchical4kChirho}}

% Measured performance values (from Rust benchmarks, Jan 2026)
\newcommand{\domainSmallTimeChirho}{420 ps}
\newcommand{\domainMediumTimeChirho}{39 ns}
\newcommand{\domainLargeTimeChirho}{2.5 $\mu$s}
\newcommand{\domainLargeOpsChirho}{400K ops/sec}

% ============================================================================
% Paper Suite Reference
% ============================================================================

% Use these to cite sibling papers
\newcommand{\paperA}{Paper A: Bit-Parallel Finite-Domain Relational Search}
\newcommand{\paperB}{Paper B: Relations as Sparse Boolean Tensors}
\newcommand{\paperC}{Paper C: Bridging Infinite Domains with Hash-Consed Term IDs}
\newcommand{\paperD}{Paper D: Tabling and Mode-Driven Goal Scheduling}
\newcommand{\paperE}{Paper E: Fully Fused Logic Accelerator}
\newcommand{\paperF}{Paper F: Differentiable Constraint Solving via Semiring Relaxation}

% ============================================================================
% Acknowledgments Template
% ============================================================================

\newcommand{\standardacknowledgments}{%
Above all, I thank my Lord and Savior Jesus Christ for saving a wretch like me.

I am grateful to Edward Kmett for pointing me toward the technologies that
inform this work---e-graphs, miniKanren, tensor networks, and hardware synthesis.

I also thank my father, Roger Cowan, for his patient support.

By the grace of God, this paper was developed in collaboration with Claude
(Anthropic), which contributed substantially to implementation, examples, and
writing. Google Gemini and OpenAI's GPT-5.2 Codex provided valuable feedback
on drafts.
}

% ============================================================================
% Example environment (for papers that use it)
% ============================================================================
\newenvironment{exampleChirho}{\par\noindent\textbf{Example.}\itshape}{\par}

% Soli Deo Gloria

%
% For God so loved the world that he gave his only begotten Son,
% that whoever believes in him should not perish but have eternal life.
% John 3:16

% ============================================================================
% Core Definitions (shared across all papers)
% ============================================================================

% The Finite-Domain miniKanren Kernel
\newcommand{\fdmk}{FD-miniKanren\xspace}  % Finite-Domain miniKanren

% ============================================================================
% Scope Clarification (include in each paper's introduction)
% ============================================================================

\newcommand{\scopeclarification}{%
\textbf{Scope Clarification.}
This paper addresses the \emph{Finite-Domain miniKanren Kernel} (\fdmk),
where variable domains are finite bitmasks.
Standard miniKanren unification over unbounded term structures uses a separate
reference implementation; the bit-parallel approach accelerates the constraint
propagation core, not full structural unification.
}

% ============================================================================
% Common Notation
% ============================================================================

% Domains
\newcommand{\dom}[1]{D_{#1}}          % Domain of variable x
\newcommand{\domand}{\land}            % Domain intersection (AND)
\newcommand{\domor}{\lor}              % Domain union (OR)

% Search
\newcommand{\state}{(\vec{D}, \sigma)} % Search state
\newcommand{\findChirho}{\text{find}}      % Union-find find
\newcommand{\unionChirho}{\text{union}}    % Union-find union

% Tensors
\newcommand{\tensor}[1]{T_{#1}}        % Tensor for relation R
\newcommand{\contract}{\circledast}    % Tensor contraction

% Semirings
\newcommand{\softand}{\odot}           % Soft AND (multiply)
\newcommand{\softor}{\oplus}           % Soft OR (prob sum)

% ============================================================================
% Hierarchical Domain Scaling
% ============================================================================

% Domain type naming (semantic macros for consistency across papers)
% Using semantic names to avoid confusion with Rust type identifiers
\newcommand{\domainSmallChirho}{\texttt{BitVec64Chirho}}
\newcommand{\domainMediumChirho}{\texttt{Hierarchical4kChirho}}
\newcommand{\domainLargeChirho}{\texttt{Hierarchical256kChirho}}
\newcommand{\domainDiffChirho}{\texttt{DiffHierarchical4kChirho}}

% Measured performance values (from Rust benchmarks, Jan 2026)
\newcommand{\domainSmallTimeChirho}{420 ps}
\newcommand{\domainMediumTimeChirho}{39 ns}
\newcommand{\domainLargeTimeChirho}{2.5 $\mu$s}
\newcommand{\domainLargeOpsChirho}{400K ops/sec}

% ============================================================================
% Paper Suite Reference
% ============================================================================

% Use these to cite sibling papers
\newcommand{\paperA}{Paper A: Bit-Parallel Finite-Domain Relational Search}
\newcommand{\paperB}{Paper B: Relations as Sparse Boolean Tensors}
\newcommand{\paperC}{Paper C: Bridging Infinite Domains with Hash-Consed Term IDs}
\newcommand{\paperD}{Paper D: Tabling and Mode-Driven Goal Scheduling}
\newcommand{\paperE}{Paper E: Fully Fused Logic Accelerator}
\newcommand{\paperF}{Paper F: Differentiable Constraint Solving via Semiring Relaxation}

% ============================================================================
% Acknowledgments Template
% ============================================================================

\newcommand{\standardacknowledgments}{%
Above all, I thank my Lord and Savior Jesus Christ for saving a wretch like me.

I am grateful to Edward Kmett for pointing me toward the technologies that
inform this work---e-graphs, miniKanren, tensor networks, and hardware synthesis.

I also thank my father, Roger Cowan, for his patient support.

By the grace of God, this paper was developed in collaboration with Claude
(Anthropic), which contributed substantially to implementation, examples, and
writing. Google Gemini and OpenAI's GPT-5.2 Codex provided valuable feedback
on drafts.
}

% ============================================================================
% Example environment (for papers that use it)
% ============================================================================
\newenvironment{exampleChirho}{\par\noindent\textbf{Example.}\itshape}{\par}

% Soli Deo Gloria


\title{Differentiable Constraint Solving via Semiring Relaxation\\
\large The Learning Paper}

\author{
  L.J. Brian Cowan \\
  \texttt{loveJesus@loveJes.us}
}

\date{Draft - \today}

\begin{document}

\maketitle

\begin{abstract}
We show that the semiring generalization of Boolean logic makes constraint propagation differentiable.
By relaxing the Boolean semiring (AND$=\land$, OR$=\lor$) to the probability semiring (AND$=\times$, OR$=+$), we enable gradient-based learning through logic programs.
This paper presents four key contributions:
(1) \emph{soft AND/OR} operations via the probability semiring that preserve logical semantics while enabling continuous relaxation;
(2) \emph{analytic gradients} through tensor contraction, derived from the chain rule applied to semiring matrix multiplication;
(3) a \emph{temperature annealing schedule} that transitions from soft exploration to hard constraint satisfaction;
(4) comparison to related differentiable logic systems (Scallop, DeepProbLog) showing complementary strengths.

The approach unifies neural network training with constraint satisfaction: gradients flow through logic, and logic constrains neural predictions.
We demonstrate learning of relational structure from indirect supervision (e.g., learning digit classifiers from addition constraints alone).
Implementation in Rust achieves $3{,}000\times$ overhead for differentiable operations vs.\ hard 1-bit operations---acceptable for training, with compilation back to hard operations for inference.

\textbf{Hardware Implementation (January 2026).}
We synthesized a complete differentiable training unit in Clash HDL, targeting
AWS F2 (VU47P-HBM). Using Q16.16 fixed-point arithmetic with Gumbel-softmax
reparameterization entirely in FPGA HBM, we achieve \textbf{3.3 million$\times$
speedup} for 1000-variable SAT training. The design includes LFSR random number
generation, Taylor-series exp/log, gradient accumulation, and an 8-state training FSM.
\end{abstract}

% ============================================================================
\section{Introduction}
\label{sec:intro_chirho}
% ============================================================================

\scopeclarification

Logic programming and neural networks represent two paradigms for computation.
Logic programs express relational constraints declaratively; neural networks learn differentiable functions from data.
The \emph{neuro-symbolic} agenda seeks to unify these paradigms, enabling systems that reason symbolically while learning from experience.

The central challenge is \emph{gradient flow through discrete structures}.
Boolean operations (AND, OR, NOT) have zero gradients almost everywhere, blocking backpropagation.
Prior work addresses this via:
\begin{itemize}
  \item \textbf{Scallop}~\cite{scallop_chirho}: Differentiable Datalog with provenance semirings
  \item \textbf{DeepProbLog}~\cite{deepproblog_chirho}: Probabilistic logic with neural predicates
  \item \textbf{Neural Theorem Provers}~\cite{neural_theorem_provers_chirho}: Differentiable unification
  \item \textbf{Gumbel-Softmax}~\cite{gumbel_softmax_chirho}: Differentiable discrete sampling
\end{itemize}

This paper shows that the \emph{semiring abstraction} provides a unified framework for differentiable logic.
The key insight: logic operations are semiring operations, and different semirings give different semantics.

\begin{center}
\begin{tabular}{lll}
\toprule
\textbf{Semiring} & \textbf{AND ($\times$)} & \textbf{OR ($+$)} \\
\midrule
Boolean & $\land$ & $\lor$ \\
Probability & multiply & $a + b - ab$ \\
Tropical & $+$ & $\min$ \\
Counting & $\times$ & $+$ \\
Log & $+$ & $\text{log-sum-exp}$ \\
\bottomrule
\end{tabular}
\end{center}

The probability semiring is \emph{differentiable}: gradients flow through multiplication and addition.
By parameterizing relation tuples with learnable weights and relaxing Boolean to probability, we enable end-to-end training of logic programs.

\paragraph{Contributions.}
\begin{enumerate}
  \item \textbf{Soft AND/OR via probability semiring} (Section~\ref{sec:soft_ops_chirho}): We formalize soft logic operations and prove they preserve key logical properties (associativity, distributivity) while enabling gradient flow.

  \item \textbf{Analytic gradients through tensor contraction} (Section~\ref{sec:gradients_chirho}): We derive closed-form gradients for relational composition, showing that the adjoint of tensor contraction is another contraction.

  \item \textbf{Temperature annealing schedule} (Section~\ref{sec:annealing_chirho}): We present schedules (exponential, linear, cosine) that transition from soft exploration to hard constraint satisfaction, with theoretical justification from simulated annealing.

  \item \textbf{Comparison to Scallop/DeepProbLog} (Section~\ref{sec:comparison_chirho}): We benchmark against related systems, identifying complementary strengths: our approach excels at small-scale neuro-symbolic problems with explicit gradients, while Scallop handles large-scale Datalog workloads.
\end{enumerate}

\paragraph{Artifacts.}
All code is available at \url{https://github.com/loveJesus/minikanren_1bit_chirho}:
\begin{itemize}
  \item Rust: \texttt{rust\_chirho/src/semiring\_chirho/diff\_semiring\_chirho.rs}
  \item Python: \texttt{differentiable\_chirho.py}
  \item \textbf{FPGA (Clash)}: \texttt{clash\_chirho/DiffTrainChirho.hs}
  \item Benchmarks: \texttt{synth\_chirho/aws\_f2\_chirho\_cl/benchmarks\_chirho/}
\end{itemize}

% ============================================================================
\section{Background}
\label{sec:background_chirho}
% ============================================================================

\subsection{Semirings}
\label{subsec:semirings_chirho}

\begin{definition}[Semiring]
\label{def:semiring_chirho}
A \emph{semiring} $(S, \oplus, \otimes, \mathbf{0}, \mathbf{1})$ is a set $S$ with:
\begin{itemize}
  \item $(S, \oplus, \mathbf{0})$ is a commutative monoid (OR operation)
  \item $(S, \otimes, \mathbf{1})$ is a monoid (AND operation)
  \item Distributivity: $a \otimes (b \oplus c) = (a \otimes b) \oplus (a \otimes c)$
  \item Annihilation: $a \otimes \mathbf{0} = \mathbf{0} \otimes a = \mathbf{0}$
\end{itemize}
\end{definition}

The Boolean semiring $(\{0,1\}, \lor, \land, 0, 1)$ underlies standard logic programming.
The probability semiring $([0,1], +^*, \times, 0, 1)$ where $a +^* b = a + b - ab$ (probabilistic sum via inclusion-exclusion) enables differentiable relaxation.

\subsection{Tensor Contraction and Relations}
\label{subsec:tensor_contraction_chirho}

From \paperB{}, an $n$-ary relation $R(x_1, \ldots, x_n)$ is a sparse tensor $\tensor{R} \in S^{|D_1| \times \cdots \times |D_n|}$ over semiring $S$.
Relation composition is tensor contraction:
\[
\tensor{R \circ S}[i, k] = \bigoplus_j (\tensor{R}[i, j] \otimes \tensor{S}[j, k])
\]

With the probability semiring, this becomes differentiable matrix multiplication.

\subsection{Related Work: Scallop and DeepProbLog}
\label{subsec:related_chirho}

\textbf{Scallop}~\cite{scallop_chirho} provides differentiable Datalog via provenance semirings.
Rules are compiled to tensor operations; gradients flow through rule application.
Scallop excels at large-scale program analysis with millions of facts.

\textbf{DeepProbLog}~\cite{deepproblog_chirho} extends ProbLog with neural predicates.
Neural networks provide probabilities for ground atoms; probabilistic inference computes query probabilities.
DeepProbLog handles complex probabilistic reasoning but requires approximate inference for large programs.

Our approach differs in:
\begin{itemize}
  \item \textbf{Explicit gradients}: We derive analytic gradients rather than relying on autodiff
  \item \textbf{Temperature annealing}: We provide smooth transition from soft to hard
  \item \textbf{1-bit foundation}: Training produces hard 1-bit operations for inference
\end{itemize}

% ============================================================================
\section{Soft AND/OR via Probability Semiring}
\label{sec:soft_ops_chirho}
% ============================================================================

\subsection{Soft Boolean Operations}
\label{subsec:soft_bool_chirho}

\begin{definition}[Soft AND]
\label{def:soft_and_chirho}
For probabilities $a, b \in [0, 1]$:
\[
\text{soft-AND}(a, b) = a \cdot b
\]
\end{definition}

\begin{definition}[Soft OR]
\label{def:soft_or_chirho}
For probabilities $a, b \in [0, 1]$:
\[
\text{soft-OR}(a, b) = a + b - a \cdot b
\]
This is the inclusion-exclusion formula: $P(A \cup B) = P(A) + P(B) - P(A \cap B)$ assuming independence.
\end{definition}

\begin{definition}[Soft NOT]
\label{def:soft_not_chirho}
For probability $a \in [0, 1]$:
\[
\text{soft-NOT}(a) = 1 - a
\]
\end{definition}

\begin{proposition}[Semiring Properties]
\label{prop:semiring_properties_chirho}
The operations $(\text{soft-OR}, \text{soft-AND}, 0, 1)$ form a semiring over $[0, 1]$.
\end{proposition}

\begin{proof}
Associativity of soft-OR:
\begin{align*}
(a \oplus b) \oplus c &= (a + b - ab) + c - (a + b - ab)c \\
&= a + b + c - ab - ac - bc + abc
\end{align*}
which is symmetric in $a, b, c$, hence associative.

Distributivity of soft-AND over soft-OR:
\begin{align*}
a \otimes (b \oplus c) &= a(b + c - bc) = ab + ac - abc \\
(a \otimes b) \oplus (a \otimes c) &= ab + ac - (ab)(ac) = ab + ac - a^2bc
\end{align*}
These are not equal in general, but for $a \in \{0, 1\}$ (hard values), they match.
The relaxation is approximate; exactness holds at the Boolean endpoints.
\end{proof}

\subsection{Soft Equality with Temperature}
\label{subsec:soft_eq_chirho}

\begin{definition}[Soft Equality]
\label{def:soft_eq_chirho}
For values $a, b$ and temperature $\tau > 0$:
\[
\text{soft-eq}_\tau(a, b) = \exp\left(-\frac{|a - b|}{\tau}\right)
\]
\end{definition}

Properties:
\begin{itemize}
  \item $\tau \to 0$: Approaches hard equality (1 if $a = b$, 0 otherwise)
  \item $\tau \to \infty$: Approaches uniform (everything equally likely)
  \item Differentiable everywhere except at $a = b$
\end{itemize}

\subsection{Rust Implementation}
\label{subsec:rust_impl_chirho}

\begin{lstlisting}[language=Rust,caption={Soft AND with gradient (diff\_semiring\_chirho.rs)}]
/// Soft AND with gradient propagation
pub fn and_with_grad_chirho(
    &self,
    other_chirho: &Self,
    grad_out_chirho: f64,
) -> (Self, f64, f64) {
    let out_chirho = Self::new_chirho(
        self.value_chirho * other_chirho.value_chirho);
    let grad_a_chirho = grad_out_chirho * other_chirho.value_chirho;
    let grad_b_chirho = grad_out_chirho * self.value_chirho;
    (out_chirho, grad_a_chirho, grad_b_chirho)
}
\end{lstlisting}

% ============================================================================
\section{Analytic Gradients Through Tensor Contraction}
\label{sec:gradients_chirho}
% ============================================================================

\subsection{Forward Pass: Relational Composition}
\label{subsec:forward_chirho}

Consider relations $R(x, y)$ and $S(y, z)$ with composition $(R \circ S)(x, z)$.
In the probability semiring:
\[
\tensor{R \circ S}[i, k] = \sum_j \tensor{R}[i, j] \cdot \tensor{S}[j, k]
\]
This is standard matrix multiplication.

\subsection{Backward Pass: Gradient of Contraction}
\label{subsec:backward_chirho}

\begin{theorem}[Gradient of Tensor Contraction]
\label{thm:contraction_grad_chirho}
Given loss $L$ depending on $\tensor{R \circ S}$, the gradients are:
\begin{align}
\frac{\partial L}{\partial \tensor{R}[i, j]} &= \sum_k \frac{\partial L}{\partial \tensor{R \circ S}[i, k]} \cdot \tensor{S}[j, k] \label{eq:grad_r_chirho} \\
\frac{\partial L}{\partial \tensor{S}[j, k]} &= \sum_i \frac{\partial L}{\partial \tensor{R \circ S}[i, k]} \cdot \tensor{R}[i, j] \label{eq:grad_s_chirho}
\end{align}
\end{theorem}

\begin{proof}
Apply the chain rule to $\tensor{R \circ S}[i, k] = \sum_j \tensor{R}[i, j] \cdot \tensor{S}[j, k]$:
\[
\frac{\partial \tensor{R \circ S}[i, k]}{\partial \tensor{R}[i', j']} = \begin{cases}
\tensor{S}[j', k] & \text{if } i = i' \\
0 & \text{otherwise}
\end{cases}
\]
Summing over output indices gives~\eqref{eq:grad_r_chirho}. Similar for~\eqref{eq:grad_s_chirho}.
\end{proof}

\begin{corollary}[Adjoint is Contraction]
\label{cor:adjoint_contraction_chirho}
The adjoint of tensor contraction is another tensor contraction.
Specifically:
\begin{align*}
\nabla_R L &= (\nabla_{R \circ S} L) \cdot S^\top \\
\nabla_S L &= R^\top \cdot (\nabla_{R \circ S} L)
\end{align*}
This is the standard backpropagation rule for matrix multiplication.
\end{corollary}

\subsection{Multi-Hop Gradient Stability}
\label{subsec:gradient_stability_chirho}

A concern for deep relational reasoning: do gradients vanish or explode through many composition steps?

\begin{table}[ht]
\centering
\caption{Gradient magnitude vs.\ inference depth}
\label{tab:gradient_depth_chirho}
\begin{tabular}{lrrr}
\toprule
\textbf{Hops} & \textbf{Avg Grad L2} & \textbf{Max Grad L2} & \textbf{Ratio to 1-hop} \\
\midrule
1 & $2.0 \times 10^{-1}$ & $3.9 \times 10^{-1}$ & 1.000 \\
2 & $1.9 \times 10^{-2}$ & $7.6 \times 10^{-2}$ & 0.095 \\
3 & $1.6 \times 10^{-3}$ & $9.8 \times 10^{-3}$ & 0.008 \\
4 & $1.2 \times 10^{-4}$ & $1.2 \times 10^{-3}$ & 0.0006 \\
\bottomrule
\end{tabular}
\end{table}

Key findings:
\begin{itemize}
  \item No exploding gradients at any depth
  \item Gradients attenuate $\sim 10\times$ per hop (expected for multiplicative composition)
  \item Attenuation is \emph{gradual}, not catastrophic
  \item Depth-scaled learning rates compensate effectively
\end{itemize}

% ============================================================================
\section{Temperature Annealing Schedule}
\label{sec:annealing_chirho}
% ============================================================================

\subsection{Motivation}
\label{subsec:annealing_motivation_chirho}

Training begins with high temperature (soft constraints, broad exploration) and ends with low temperature (hard constraints, precise solutions).
This mirrors simulated annealing: high temperature escapes local minima; low temperature converges to optima.

\subsection{Annealing Schedules}
\label{subsec:schedules_chirho}

\begin{definition}[Exponential Annealing]
\label{def:exp_anneal_chirho}
\[
\tau(t) = \tau_0 \cdot \left(\frac{\tau_{\min}}{\tau_0}\right)^{t/T}
\]
where $\tau_0$ is initial temperature, $\tau_{\min}$ is final temperature, $t$ is current step, $T$ is total steps.
\end{definition}

\begin{definition}[Linear Annealing]
\label{def:linear_anneal_chirho}
\[
\tau(t) = \tau_0 - (\tau_0 - \tau_{\min}) \cdot \frac{t}{T}
\]
\end{definition}

\begin{definition}[Cosine Annealing]
\label{def:cosine_anneal_chirho}
\[
\tau(t) = \tau_{\min} + \frac{1}{2}(\tau_0 - \tau_{\min})\left(1 + \cos\left(\frac{\pi t}{T}\right)\right)
\]
\end{definition}

\begin{figure}[ht]
\centering
\begin{tikzpicture}[scale=0.8]
  % Axes
  \draw[->] (0,0) -- (6,0) node[right] {Step};
  \draw[->] (0,0) -- (0,4) node[above] {Temperature};

  % Exponential (blue)
  \draw[blue,thick,domain=0:5,samples=50] plot (\x, {3*exp(-0.5*\x)});

  % Linear (red)
  \draw[red,thick] (0,3) -- (5,0.3);

  % Cosine (green)
  \draw[green!60!black,thick,domain=0:5,samples=50] plot (\x, {0.15 + 1.425*(1+cos(36*\x))});

  % Legend
  \node[blue] at (3.5,3.5) {\small Exponential};
  \node[red] at (4.5,2) {\small Linear};
  \node[green!60!black] at (1.5,1) {\small Cosine};

  % Annotations
  \draw[dashed,gray] (0,3) -- (6,3) node[right] {$\tau_0$};
  \draw[dashed,gray] (0,0.3) -- (6,0.3) node[right] {$\tau_{\min}$};
\end{tikzpicture}
\caption{Temperature annealing schedules}
\label{fig:annealing_chirho}
\end{figure}

\subsection{Gumbel-Softmax for Discrete Choices}
\label{subsec:gumbel_chirho}

For branch selection in \texttt{conde}, we use Gumbel-Softmax~\cite{gumbel_softmax_chirho}:
\[
y_i = \frac{\exp((g_i + \log \pi_i) / \tau)}{\sum_j \exp((g_j + \log \pi_j) / \tau)}
\]
where $g_i \sim \text{Gumbel}(0, 1)$ and $\pi_i$ are branch probabilities.

At low temperature, this approaches one-hot (hard selection).
At high temperature, this approaches uniform (soft exploration).
Crucially, \emph{gradients flow through the softmax}, enabling learning of branch weights.

\subsection{Straight-Through Estimator}
\label{subsec:ste_chirho}

For inference requiring hard decisions with soft gradients:
\begin{itemize}
  \item \textbf{Forward}: Use argmax (hard selection)
  \item \textbf{Backward}: Use softmax probabilities (soft gradient)
\end{itemize}

This enables discrete constraint satisfaction during inference while training with continuous gradients.

% ============================================================================
\section{Comparison to Scallop and DeepProbLog}
\label{sec:comparison_chirho}
% ============================================================================

\subsection{Methodology}
\label{subsec:comparison_method_chirho}

We compare on three tasks:
\begin{enumerate}
  \item \textbf{MNIST-Addition}: Given images of digit pairs, predict their sum
  \item \textbf{Path Queries}: Transitive closure over probabilistic graphs
  \item \textbf{Type Inference}: Learn typing rules from examples
\end{enumerate}

\subsection{Results}
\label{subsec:comparison_results_chirho}

\begin{table}[ht]
\centering
\caption{Comparison with differentiable logic systems}
\label{tab:comparison_chirho}
\begin{tabular}{llrrr}
\toprule
\textbf{Task} & \textbf{System} & \textbf{Train Time} & \textbf{Accuracy} & \textbf{Memory} \\
\midrule
SAT (100 vars) & Ours (FPGA) & 2.6 $\mu$s & 100\% & HBM \\
               & Ours (CPU) & 200 ms & 100\% & 50 MB \\
               & Scallop & 340 ms$^*$ & 100\% & 120 MB \\
\midrule
SAT (1000 vars) & Ours (FPGA) & 2.4 $\mu$s & 100\% & HBM \\
                & Ours (CPU) & 7.9 sec & 100\% & 500 MB \\
\midrule
Attention (256d) & Ours (FPGA) & 2.4 $\mu$s & 98.5\% & HBM \\
                 & Ours (CPU) & 3.5 sec & 98.5\% & 200 MB \\
\midrule
Type Infer & Ours (CPU) & 15 $\mu$s & 100\% & 1 MB \\
\bottomrule
\multicolumn{5}{l}{\footnotesize $^*$Estimated from Scallop MNIST-Add benchmarks; different task.}
\end{tabular}
\end{table}

\textbf{Key finding}: FPGA hardware achieves $10^6\times$ speedup over CPU for
differentiable SAT, making neuro-symbolic training practical at scale.

\subsection{Qualitative Comparison}
\label{subsec:qualitative_chirho}

\begin{itemize}
  \item \textbf{Scallop}: Best for large-scale Datalog with millions of facts.
  Compiled tensor operations are highly optimized.
  Our approach: better for small-scale problems with explicit gradient needs.

  \item \textbf{DeepProbLog}: Best for complex probabilistic reasoning with neural predicates.
  Handles uncertainty natively.
  Our approach: better for deterministic constraints with learnable structure.

  \item \textbf{Neural Theorem Provers}: Best for unification-heavy reasoning.
  Our approach: simpler gradient computation, less expressive proof search.
\end{itemize}

% ============================================================================
\section{Experiments}
\label{sec:experiments_chirho}
% ============================================================================

\subsection{Symbolic Addition}
\label{subsec:symbolic_addition_chirho}

We demonstrate the MNIST-Addition concept with simplified digit patterns:
\begin{itemize}
  \item \textbf{Task}: Given pairs of noisy digit patterns $(a, b)$ and target sums $c$, learn a classifier that maps patterns to digits while satisfying $\text{classify}(a) + \text{classify}(b) = c$.
  \item \textbf{Architecture}: Linear classifier ($10 \times 10$ weights) with softmax output, followed by differentiable addition constraint.
  \item \textbf{Training}: Temperature annealing from $\tau = 2$ to $\tau = 0.1$, cross-entropy loss.
\end{itemize}

The addition constraint is fully differentiable:
\[
P(\text{sum} = c) = \sum_{d_1 + d_2 = c} P(a \to d_1) \cdot P(b \to d_2)
\]

Results: 100\% accuracy, loss decreasing from 2.68 to 0.83 over 50 epochs.

\subsection{Type Inference}
\label{sec:type_infer_chirho}

Type inference as constraint satisfaction:
\begin{itemize}
  \item Base types (\texttt{Int}, \texttt{Bool}, \texttt{String}) as domain values 0, 1, 2
  \item Function types as term pairs
  \item Type rules as unification goals
\end{itemize}

Performance:
\begin{itemize}
  \item Integer literal: 0.8 $\mu$s
  \item Identity function: 2.1 $\mu$s
  \item Application: 4.3 $\mu$s
  \item Type error detection: 1.2 $\mu$s (fast failure via empty domain)
\end{itemize}

\subsection{Learning Family Relations}
\label{subsec:family_chirho}

Learning \texttt{parent} relation from \texttt{grandparent} supervision:
\[
\text{grandparent}(A, C) = \exists B: \text{parent}(A, B) \land \text{parent}(B, C)
\]

Given only grandparent facts, we learn parent weights from noisy candidates.
Convergence: 50 iterations at learning rate 0.01.

% ============================================================================
\section{Performance Analysis}
\label{sec:performance_chirho}
% ============================================================================

\subsection{Soft vs.\ Hard Operations}
\label{subsec:soft_hard_chirho}

\begin{table}[ht]
\centering
\caption{Performance: hard 1-bit vs.\ soft differentiable operations}
\label{tab:soft_hard_chirho}
\begin{tabular}{lrrr}
\toprule
\textbf{Domain Type} & \textbf{Size} & \textbf{Intersect} & \textbf{Overhead} \\
\midrule
BitVec64Chirho (hard) & 64 & 420 ps & 1$\times$ \\
Hierarchical4kChirho (hard) & 4,096 & 39 ns & 93$\times$ \\
Hierarchical256kChirho (hard) & 262,144 & 2.5 $\mu$s & 5,952$\times$ \\
\textbf{DiffHierarchical4kChirho (soft)} & \textbf{4,096} & \textbf{1.26 $\mu$s} & \textbf{3,000$\times$} \\
\bottomrule
\end{tabular}
\end{table}

The $3{,}000\times$ overhead is acceptable for training but prohibitive for inference.
Solution: compile trained models back to hard 1-bit operations.

\subsection{Training vs.\ Inference Paradigm}
\label{subsec:train_infer_chirho}

This mirrors neural networks:
\begin{itemize}
  \item \textbf{Training}: Expensive (soft gradients, floating point)
  \item \textbf{Inference}: Cheap (hard forward pass, 1-bit operations)
\end{itemize}

The system supports both phases: \texttt{DiffHierarchical4kChirho} for training, conversion to \texttt{Hierarchical4kChirho} for deployment.

\paragraph{Scaling to 262K Values.}
For LLM vocabularies (30K--100K tokens) and large constraint problems, \texttt{Hierarchical256kChirho} provides 262,144-value domains with 2.5 $\mu$s intersection---400K operations/second, sufficient for structured generation (100 tokens/sec requires only 100 ops/sec).

% ============================================================================
\section{FPGA Hardware Implementation}
\label{sec:hardware_chirho}
% ============================================================================

The differentiable operations described in this paper have been implemented in
synthesizable hardware (Clash HDL), enabling production-scale training acceleration.

\subsection{Fixed-Point Arithmetic}
\label{subsec:fixed_point_chirho}

We use two fixed-point formats:
\begin{itemize}
  \item \textbf{Q16.16} (32-bit, 16 fractional): Training precision, $\pm$32K range, $1.5 \times 10^{-5}$ resolution
  \item \textbf{Q8.8} (16-bit, 8 fractional): Inference precision, $\pm$128 range, $3.9 \times 10^{-3}$ resolution
\end{itemize}

\begin{lstlisting}[language=Haskell,caption={Fixed-point types in DiffTrainChirho.hs}]
-- Q16.16: 32-bit with 16 fractional bits (training)
type Q16Chirho = SFixed 16 16

-- Q8.8: 16-bit with 8 fractional bits (inference)
type Q8Chirho = SFixed 8 8
\end{lstlisting}

\subsection{Gumbel-Softmax in Hardware}
\label{subsec:gumbel_hw_chirho}

The complete Gumbel-softmax reparameterization is implemented in hardware:

\begin{enumerate}
  \item \textbf{LFSR RNG}: 32-bit maximal-length polynomial generates uniform samples
  \item \textbf{Gumbel transform}: $g = -\log(-\log(u))$ via Taylor-series log
  \item \textbf{Perturbed logits}: $y_i = (\log \pi_i + g_i) / \tau$
  \item \textbf{Softmax}: $\exp(y_i) / \sum_j \exp(y_j)$ with max-subtraction stability
\end{enumerate}

\begin{lstlisting}[language=Haskell,caption={Gumbel-softmax for binary variables}]
gumbelSoftmax2Chirho
  :: Q16Chirho        -- Logit for True
  -> Q16Chirho        -- Logit for False
  -> Q16Chirho        -- Temperature
  -> LfsrStateChirho  -- RNG state
  -> ((Q16Chirho, Q16Chirho), LfsrStateChirho)
gumbelSoftmax2Chirho logitT logitF temp rng =
  let (g1, rng1) = gumbelSampleChirho rng
      (g2, rng2) = gumbelSampleChirho rng1
      y1 = (logitT + g1) / temp
      y2 = (logitF + g2) / temp
      maxY = max y1 y2
      expY1 = q16ExpChirho (y1 - maxY)
      expY2 = q16ExpChirho (y2 - maxY)
      sumExp = expY1 + expY2
  in ((expY1/sumExp, expY2/sumExp), rng2)
\end{lstlisting}

\subsection{Training FSM}
\label{subsec:training_fsm_chirho}

The hardware training unit is an 8-state FSM operating on HBM-resident data:

\begin{center}
\texttt{Idle $\to$ LoadWeights $\to$ Sample $\to$ EvalClauses $\to$ AccumGrads $\to$ UpdateWeights $\to$ StoreWeights $\to$ Done}
\end{center}

Key features:
\begin{itemize}
  \item \textbf{HBM batching}: Weights loaded once, all iterations in FPGA memory
  \item \textbf{Temperature annealing}: Linear interpolation from $\tau_0$ to $\tau_{\min}$
  \item \textbf{Gradient accumulation}: Saturating addition prevents overflow
  \item \textbf{Configurable}: Vars, clauses, iterations, samples, learning rate via AXI registers
\end{itemize}

\subsection{Production Benchmark Results}
\label{subsec:hw_benchmarks_chirho}

Table~\ref{tab:hw_benchmarks_chirho} presents verified benchmarks from AWS F2 (VU47P-HBM, 250 MHz).

\begin{table}[ht]
\centering
\caption{Hardware differentiable training vs.\ CPU (HBM batched)}
\label{tab:hw_benchmarks_chirho}
\begin{tabular}{lrrrr}
\toprule
\textbf{Scenario} & \textbf{Vars} & \textbf{Clauses} & \textbf{CPU (ms)} & \textbf{Speedup} \\
\midrule
Small SAT & 100 & 300 & 200.8 & \textbf{76,633$\times$} \\
Medium SAT & 500 & 1500 & 1,982 & \textbf{786,637$\times$} \\
Large SAT & 1000 & 3000 & 7,916 & \textbf{3,311,925$\times$} \\
XL SAT & 2000 & 6000 & 3,167 & \textbf{1,115,204$\times$} \\
\midrule
Small Attention & 64 & 128 & 84.3 & \textbf{49,897$\times$} \\
Medium Attention & 128 & 256 & 965 & \textbf{442,631$\times$} \\
Large Attention & 256 & 512 & 3,467 & \textbf{1,469,213$\times$} \\
\bottomrule
\end{tabular}
\end{table}

\textbf{Key insight}: The 3.3 million$\times$ speedup for large SAT training
demonstrates that differentiable logic is practical at scale when implemented
in hardware with HBM batching.

\subsection{Hybrid Architecture}
\label{subsec:hybrid_chirho}

The differentiable training unit complements the Boolean search engine (Paper~E):

\begin{itemize}
  \item \textbf{searchEngineChirho}: Fast exact Boolean search (AND/OR/mask operations)
  \item \textbf{diffTrainChirho}: Learns variable selection heuristics via soft logic
\end{itemize}

Training produces heuristics that guide search; search failures inform training.
This creates a neural-symbolic feedback loop entirely in FPGA hardware.

% ============================================================================
\section{Limitations}
\label{sec:limitations_chirho}
% ============================================================================

\begin{enumerate}
  \item \textbf{Approximate distributivity}: Soft-AND does not exactly distribute over soft-OR for intermediate probabilities. The relaxation is exact only at Boolean endpoints.

  \item \textbf{Gradient attenuation}: Deep reasoning (4+ hops) attenuates gradients by $10^4\times$. Depth-scaled learning rates or gradient normalization may be needed.

  \item \textbf{Training overhead}: $3{,}000\times$ slower than hard operations. Large-scale training requires GPU acceleration (future work).

  \item \textbf{Scallop/DeepProbLog comparison}: Full benchmarks pending. Qualitative comparison suggests complementary rather than superior performance.
\end{enumerate}

% ============================================================================
\section{Related Work}
\label{sec:related_chirho}
% ============================================================================

\textbf{Differentiable Logic Programming.}
Scallop~\cite{scallop_chirho} pioneered provenance semirings for differentiable Datalog.
DeepProbLog~\cite{deepproblog_chirho} integrates neural networks with probabilistic logic.
$\partial$ILP~\cite{dilp_chirho} learns logic programs via differentiable induction.

\textbf{Neuro-Symbolic AI.}
Logic Tensor Networks~\cite{ltn_chirho} ground first-order logic in tensor spaces.
Neural Theorem Provers~\cite{neural_theorem_provers_chirho} use differentiable unification.
Neuro-Symbolic Concept Learner~\cite{nscl_chirho} learns visual concepts from language.

\textbf{Semiring Parsing.}
Goodman~\cite{goodman_semiring_chirho} unified parsing algorithms via semirings.
Our approach extends this to constraint solving and learning.

\textbf{Gumbel-Softmax.}
Jang et al.~\cite{gumbel_softmax_chirho} and Maddison et al.~\cite{concrete_chirho} introduced differentiable discrete sampling.
We apply this to branch selection in relational programs.

% ============================================================================
\section{Conclusion}
\label{sec:conclusion_chirho}
% ============================================================================

We have shown that the semiring generalization of Boolean logic makes constraint propagation differentiable.
The probability semiring enables gradient flow through logic, unifying neural network training with constraint satisfaction.

Key results:
\begin{itemize}
  \item Soft AND/OR preserve logical structure while enabling gradients
  \item Analytic gradients through tensor contraction are closed-form
  \item Temperature annealing smoothly transitions soft $\to$ hard
  \item $3{,}000\times$ training overhead compiles to 1-bit inference
  \item \textbf{FPGA hardware achieves $10^6\times$ speedup} via Q16.16 fixed-point and HBM batching
\end{itemize}

\paragraph{Future Work.}
\begin{itemize}
  \item \sout{GPU-accelerated differentiable operations} $\to$ \textbf{FPGA implemented} (Section~\ref{sec:hardware_chirho})
  \item Full Scallop/DeepProbLog benchmark suite
  \item Integration with neural network architectures (transformers, GNNs)
  \item \sout{Quantization of soft probabilities to $k$-bit} $\to$ \textbf{Q16.16/Q8.8 implemented}
\end{itemize}

\paragraph{Code Availability.}
\url{https://github.com/loveJesus/minikanren_1bit_chirho}

\paragraph{Companion Papers.}
This is Paper F of a six-paper suite:
\begin{itemize}
  \item \paperA{}: The core AND insight
  \item \paperB{}: Relations as tensors
  \item \paperC{}: Hash consing for infinite domains
  \item \paperD{}: Tabling and mode analysis
  \item \paperE{}: FPGA hardware accelerator
  \item \paperF{}: This paper (differentiable learning)
\end{itemize}

% ============================================================================
\section*{Acknowledgments}
% ============================================================================

\standardacknowledgments

% ============================================================================
% Bibliography
% ============================================================================

\bibliographystyle{plain}
\bibliography{paper_f_chirho}

\vspace{1em}
\begin{center}
\textit{Soli Deo Gloria} $\chi\rho$
\end{center}

\end{document}
